% ============================================================
% РАЗДЕЛ 2. Анализ существующих решений
% ============================================================
\section{АНАЛИЗ СУЩЕСТВУЮЩИХ РЕШЕНИЙ ДЛЯ ХРАНЕНИЯ И КЭШИРОВАНИЯ МОДЕЛЕЙ}

В данном разделе исследуются архитектуры и свойства наиболее распространённых систем распределённого хранения и кэширования, рассматриваемых в качестве кандидатных решений для задачи ускорения холодного старта ML-инференса в Kubernetes. Для каждого решения проводится анализ применимости к сценарию бессерверного инференса: определяется, в какой мере система устраняет повторные загрузки весов, обеспечивает локальность данных на worker-нодах и сохраняет прозрачный файловый интерфейс для runtime-контейнера. На основе анализа формулируется обоснование разработки специализированного решения.

\subsection{Критерии сравнения решений}
\label{sec:criteria}

Чтобы сравнение систем было обоснованным и воспроизводимым, необходимо зафиксировать критерии, непосредственно вытекающие из требований, сформулированных в разделе~1. Каждый критерий отражает конкретный аспект задачи ускорения холодного старта.

\textbf{Задержка при холодном старте} определяет, насколько система способна уменьшить время до первого доступного байта весов модели. Ключевой характеристикой здесь является поддержка ленивой загрузки (lazy loading) или чтения по требованию: runtime-контейнер должен иметь возможность стартовать и обращаться к отдельным фрагментам файла ещё до его полной передачи.

\textbf{Поддержка Kubernetes} означает наличие официального механизма развёртывания в контейнерной среде: CSI-драйвера, Helm-чарта или Kubernetes Operator. Отсутствие этих инструментов существенно увеличивает операционную сложность интеграции.

\textbf{Локальный кэш на ноде} — возможность хранить горячие данные непосредственно на дисках worker-нод вблизи compute-ресурсов, что обеспечивает минимальную latency при повторном чтении. Решения, хранящие данные только в центральном хранилище, не позволяют избавиться от сетевого хопа при каждом обращении.

\textbf{P2P-доставка} отражает способность системы передавать данные между worker-нодами без обращения к центральному источнику, тем самым распределяя нагрузку и ускоряя одновременное масштабирование большого числа реплик инференса.

\textbf{FUSE/POSIX-интерфейс} оценивает, предоставляет ли система runtime-контейнеру стандартный файловый путь, не требующий изменений в ML-фреймворке. Наличие POSIX-семантики позволяет vLLM, Triton и аналогичным инструментам работать без модификации кода.

\textbf{Отказоустойчивость} характеризует поведение системы при выходе из строя отдельных узлов: существуют ли механизм репликации данных и автоматический fallback на альтернативный источник.

\textbf{Наличие единой точки отказа (SPOF)} указывает, имеется ли в архитектуре компонент, отказ которого делает недоступными все данные. Это критически важно для сред с требованиями высокой доступности.

\textbf{Зависимость от внешнего объектного хранилища} определяет, является ли object storage обязательным звеном в пути чтения. Обязательное обращение к S3/GCS при каждом промахе кэша существенно увеличивает latency и создаёт зависимость от внешней инфраструктуры.

\textbf{Прозрачность для runtime} оценивает, насколько решение скрывает распределённую природу данных от приложения. Идеальная прозрачность означает, что ML-контейнер «видит» обычную директорию с файлами модели.

\textbf{Сложность эксплуатации} суммирует требования к компетенциям команды и инфраструктурным зависимостям. Решения с высокой сложностью неприемлемы для небольших платформ без выделенного SRE-отдела.

\textbf{Применимость к serverless cold start} является интегральным критерием: насколько хорошо решение адаптировано к сценарию эфемерных подов, работающих с крупными бинарными файлами весов, при жёстких ограничениях на время инициализации.

Перечисленные критерии используются в сравнительной таблице~\ref{tab:solutions_comparison} и в итоговом обосновании разработки специализированного решения.

\subsection{Alluxio}
\label{sec:alluxio}

Alluxio (ранее Tachyon) --- виртуальная распределённая файловая система (Virtual Distributed File System, VDFS), выступающая слоем оркестрации данных между вычислительными фреймворками и системами постоянного хранения (Amazon S3, HDFS, NFS и другими) \cite{alluxio_arch}. Система разработана в Калифорнийском университете в Беркли и нашла применение в крупномасштабных аналитических и ML-рабочих нагрузках.

\subsubsection{Архитектура}

Alluxio строится по модели Master/Worker/Client. \textit{Ведущий узел} (Leading Master) управляет глобальными метаданными: деревом inode, распределением блоков и ёмкостью рабочих узлов. Все транзакции файловой системы записываются в распределённый журнал для обеспечения восстановления состояния при сбое. Для обеспечения высокой доступности (HA) рядом с ведущим мастером работают несколько резервных мастеров, непрерывно синхронизирующихся через журнал и готовых принять управление при отказе основного.

\textit{Worker-узлы} управляют локальными ресурсами хранения (RAM, SSD, HDD), выделенными под кэш Alluxio. Данные хранятся в виде блоков. При заполнении ёмкости применяются политики вытеснения (LRU и аналоги). Worker-узлы также выполняют операции с базовым хранилищем (Under File System, UFS): при первом обращении к данным они кэшируют блоки локально, снижая нагрузку на UFS при последующих чтениях.

\textit{Client} предоставляет приложению шлюз к кластеру Alluxio. Он взаимодействует с Master для операций с метаданными и с Worker для чтения и записи данных. Поддерживаются нативный Java API, REST, HDFS API и S3-совместимый API. Важно отметить, что клиент Alluxio никогда не обращается напрямую к базовому хранилищу: данные всегда передаются через Worker-узлы \cite{alluxio_arch}.

Интеграция с Kubernetes реализована через Helm-чарты и Alluxio Operator, который автоматизирует развёртывание, настройку и управление жизненным циклом кластера Alluxio поверх Kubernetes \cite{alluxio_k8s}. Alluxio поддерживает монтирование через FUSE (Alluxio POSIX API), что обеспечивает прозрачный доступ для приложений через стандартный файловый интерфейс.

В Enterprise-версии архитектура мастера эволюционировала в сторону децентрализации: владение метаданными распределяется между Worker-узлами через согласованное хеширование, что устраняет централизованный мастер как узкое место и позволяет обходить его при чтении данных.

\subsubsection{Преимущества}

Alluxio обеспечивает высокую производительность за счёт приоритетного кэширования в оперативной памяти (memory-first tiered architecture). Унифицированное пространство имён (Unified Namespace) позволяет монтировать разнородные хранилища в единую логическую файловую систему, что актуально для платформ с несколькими объектными хранилищами. Поддержка локальности данных выгодна для ML-нагрузок: когда модельные данные кэшируются на ноде, где запущен инференс, задержка чтения снижается до уровня локального диска.

\subsubsection{Недостатки}

Первый существенный недостаток --- \textit{сложность эксплуатации}. Даже при наличии Kubernetes Operator настройка HA-конфигурации мастеров, управление eviction-политиками, диагностика утечек памяти в Worker-узлах и оптимизация распределения блоков требуют глубоких знаний архитектуры системы. Операционная нагрузка несопоставима с задачей кэширования весов модели.

Второй недостаток --- \textit{высокое потребление памяти}. Alluxio является memory-centric системой: для достижения декларируемой производительности Worker-узлы должны располагать значительным объёмом оперативной памяти. При кэшировании нескольких больших моделей (десятки гигабайт каждая) потребности в памяти конкурируют с самим GPU-инференсом.

Третий недостаток --- \textit{модель лицензирования}. Ключевые функции для производственных ML-сред: децентрализованная архитектура, расширенные политики кэширования и продвинутые AI/ML-оптимизации --- доступны только в Enterprise Edition \cite{alluxio_pricing}. Стоимость не публикуется и определяется индивидуально, что создаёт финансовую неопределённость.

Четвёртый недостаток --- \textit{несовместимость с serverless-подами}. Клиент Alluxio предполагает наличие состояния и устойчивого соединения с кластером. В эфемерных подах Kubernetes каждая инициализация клиента добавляет задержку к холодному старту, а данные, кэшированные по запросу короткоживущего пода, могут быть вытеснены прежде, чем повторно использованы.

\subsubsection{Вывод по Alluxio}

Alluxio технически близок к задаче: он обеспечивает кэширование данных с локальностью и поддерживает FUSE-монтирование. Однако он представляет собой тяжёлый data orchestration layer общего назначения, предназначенный для аналитики и распределённого ML-обучения. Для узкой задачи chunk-кэширования весов модели в serverless-инференсе это архитектурно и операционно избыточное решение.

\subsection{JuiceFS}
\label{sec:juicefs}

JuiceFS --- высокопроизводительная распределённая файловая система с POSIX-совместимым интерфейсом, ориентированная на облачные среды \cite{juicefs_intro}. Community Edition распространяется под лицензией Apache~2.0. Система нашла применение в ML-платформах благодаря нативной интеграции с Kubernetes и возможности использовать уже существующие объектные хранилища как backend.

\subsubsection{Архитектура}

Фундаментальной особенностью JuiceFS является \textit{строгое разделение данных и метаданных}. Архитектура включает три компонента.

\textit{Metadata Engine} хранит всю информацию о файловой системе: имена файлов, права доступа, временные метки, структуру директорий и, что критически важно, отображение файлов на физические блоки данных. Community Edition поддерживает несколько бэкендов: Redis, TiKV, PostgreSQL, MySQL/MariaDB \cite{juicefs_design}. Выбор Redis обеспечивает минимальные задержки (порядка миллисекунд) за счёт хранения в оперативной памяти, TiKV обеспечивает горизонтальную масштабируемость для платформ с миллиардами объектов.

\textit{Object Storage} хранит содержимое файлов в виде блоков. При записи файл логически делится на чанки по 64~МБ, каждый из которых при сбросе (flush) разбивается на блоки по 4~МБ, параллельно загружаемые в объектное хранилище \cite{juicefs_design}. JuiceFS поддерживает Amazon S3, MinIO, Google Cloud Storage и другие S3-совместимые системы. В самом объектном хранилище исходные файлы не представлены: там хранятся пронумерованные блоки, а их связь с файловой иерархией поддерживается через Metadata Engine.

\textit{Клиент JuiceFS} реализует файловый интерфейс (FUSE, CSI-драйвер, S3 Gateway) и взаимодействует одновременно с Metadata Engine (для получения структуры файла и местоположения блоков) и с объектным хранилищем (для чтения и записи данных).

В Kubernetes JuiceFS развёртывается через CSI-драйвер, который реализует архитектуру с выделенным Mount Pod на каждом узле. Этот под исполняет клиент FUSE и обслуживает все application pods на узле, использующие один и тот же PVC. Такое разделение позволяет перезапускать CSI-компоненты без прерывания доступа к данным для уже запущенных приложений \cite{juicefs_perf}.

Кэширование на локальных дисках является ключевым механизмом ускорения: при повторном чтении одних и тех же блоков JuiceFS обслуживает запросы с локального NVMe-диска вместо обращения к объектному хранилищу. По данным разработчиков, при последовательном чтении крупных файлов с включённым локальным кэшем пропускная способность может вырасти с~4{,}7~ГБ/с до~13{,}2~ГБ/с, а задержка случайного чтения --- снизиться с~29~мс до~0{,}3~мс \cite{juicefs_perf}.

\subsubsection{Ограничения для serverless cold start}

\textit{Зависимость от объектного хранилища.} JuiceFS не имеет собственного слоя хранения данных и обязательно требует внешнего S3"=совместимого хранилища. При промахе кэша (cache miss) клиент обращается к объектному хранилищу, задержки которого на порядок превышают задержки локального диска.

\textit{Metadata Engine как потенциальная точка отказа.} При использовании Redis в качестве бэкенда метаданных и отсутствии HA-конфигурации (Redis Sentinel или Redis Cluster) отказ экземпляра Redis делает недоступной всю файловую систему, даже если объектное хранилище функционирует нормально. Настройка кластерного Redis добавляет операционную сложность.

\textit{Проблема прогрева кэша в эфемерных средах.} В serverless-сценарии поды часто запускаются на «чистых» узлах, где локальный кэш пуст. В отличие от долгоживущих приложений, эфемерный под не успевает накопить полезный кэш до своего завершения. При масштабировании на новые узлы все читаемые блоки будут приходить из объектного хранилища --- именно в тот момент, когда критична latency первого байта.

\textit{Накладные расходы FUSE.} Реализация файловой системы в пространстве пользователя через FUSE вносит дополнительные переключения контекста при каждом файловом вызове. При большом числе мелких операций чтения это снижает эффективную пропускную способность одного потока.

\textit{Шторм монтирований при scale-out.} При одновременном масштабировании на большое число новых узлов создание Mount Pod'ов для каждого узла увеличивает нагрузку на API-сервер Kubernetes и может задерживать инициализацию новых подов.

\subsubsection{Вывод по JuiceFS}

JuiceFS представляет ценный архитектурный ориентир: разделение данных и метаданных, использование Redis как быстрого metadata backend, FUSE-интерфейс --- все эти идеи применены в разрабатываемой системе. Однако JuiceFS является файловой системой общего назначения с обязательной зависимостью от объектного хранилища, тогда как для решения задачи требуется специализированный cache/fetch слой с P2P-передачей чанков между нодами кластера.

\subsection{Dragonfly, CubeFS и Rook/Ceph}
\label{sec:alternatives}

\subsubsection{Dragonfly}

Dragonfly --- P2P-система распределения файлов и контейнерных образов, разработанная Alibaba Cloud и переданная в Cloud Native Computing Foundation (CNCF) \cite{dragonfly_blog}. Система ориентирована на сценарии эффективного распространения крупных файлов на большое число узлов одновременно, что делает её концептуально близкой к задаче distributing весов ML-моделей.

\predheading{Архитектура.} Dragonfly сочетает элементы CDN и P2P. \textit{Scheduler} является центральным компонентом: он отслеживает доступность узлов-пиров, их загрузку и местоположение данных, определяя оптимальные источники для каждого запроса. \textit{Seed Peer} --- специализированный пир, кэширующий данные и выступающий первичным источником для других пиров, особенно когда исходный сервер недоступен или перегружен. \textit{Peer (dfdaemon)} работает на каждом узле Kubernetes как DaemonSet, перехватывает запросы на загрузку образов и файлов, перенаправляя их в P2P-сеть. При загрузке каждый пир одновременно становится источником для других, распределяя нагрузку.

\predheading{Применимость к ML-инференсу.} Dragonfly хорошо зарекомендовал себя для параллельной дистрибуции образов контейнеров на сотни нод. Та же модель применима к весам ML-моделей: при первой загрузке модели на один из узлов она кэшируется и становится доступной для скачивания другими узлами через P2P-сеть вместо обращения к центральному репозиторию. Это снижает нагрузку на Hugging Face или внутренний model registry и ускоряет масштабирование.

\predheading{Ограничения.} Dragonfly не предоставляет стандартного FUSE/POSIX-интерфейса: система занимается доставкой файлов, но не монтированием файловой системы. Для предоставления ML-runtime доступа к весам модели через обычный файловый путь потребуется дополнительный уровень. Scheduler и Seed Peer требуют отдельного развёртывания и поддержки, увеличивая операционную нагрузку. Наличие централизованного Scheduler вводит компонент, отказ которого нарушает координацию P2P-сети, хотя фактическая доставка данных продолжается между пирами. Тем не менее Dragonfly является ближайшим по духу предшественником P2P-подхода, разработанного в данной работе: идея обмена фрагментами данных между узлами кластера вместо многократного скачивания из центрального источника лежит в основе обоих решений.

\subsubsection{CubeFS}

CubeFS --- облачно-ориентированная распределённая файловая система уровня CNCF Graduated \cite{cubefs_arch}, разработанная JD.com. Система предоставляет унифицированное хранилище с поддержкой нескольких протоколов доступа и предназначена для высоконагруженных производственных сред.

\predheading{Архитектура.} CubeFS включает несколько подсистем. \textit{Master-узлы} управляют ресурсами кластера: отслеживают состояние узлов метаданных и узлов данных, управляют томами и обеспечивают согласованность конфигурации с использованием Raft и RocksDB. \textit{Узлы метаданных} управляют шардами метаданных (metadata partition). Каждый шард содержит две B-Tree структуры в памяти: inode B-Tree и dentry B-Tree. Минимальная конфигурация требует трёх экземпляров. \textit{Data Nodes} управляют шардами данных (Data Partition) и образуют репликационные группы. CubeFS поддерживает как репликацию данных, так и Erasure Coding (через подсистему Blobstore) для более эффективного использования дискового пространства. Объектный шлюз (Object Subsystem) предоставляет S3-совместимый API.

\predheading{Поддерживаемые протоколы.} CubeFS поддерживает POSIX, S3 и HDFS-интерфейсы, а также CSI-драйвер для Kubernetes. Мультипротокольная поддержка делает систему гибкой для смешанных рабочих нагрузок, где различные компоненты ML-платформы используют разные интерфейсы доступа к данным.

\predheading{Применимость к ML-инференсу.} CubeFS обеспечивает высокую пропускную способность для крупных файлов за счёт параллелизма Data Nodes. POSIX-совместимость позволяет использовать систему для хранения и одновременного доступа к файлам весов модели из нескольких реплик инференса. Репликация данных обеспечивает отказоустойчивость при выходе из строя отдельных узлов.

\predheading{Ограничения.} Сложность развёртывания CubeFS сопоставима с Ceph: кластер требует управления несколькими типами узлов (Master, узел метаданных, DataNode), настройки репликации, мониторинга состояния и обеспечения кворума узлов метаданных. Для небольших и средних ML-платформ операционные требования несоразмерны задаче. Кроме того, CubeFS проектировалась как система хранения общего назначения, а не как специализированный кэш для бессерверного инференса: в ней отсутствуют P2P-доставка чанков и встроенный механизм lazy loading для ML-весов.

\subsubsection{Rook/Ceph}

Rook --- cloud-native оркестратор для Kubernetes, автоматизирующий развёртывание и управление распределёнными системами хранения \cite{rook_docs}. В наиболее распространённой конфигурации Rook используется для оркестрации Ceph --- одной из наиболее функционально богатых open-source-систем хранения данных.

\predheading{Архитектура Ceph.} Ceph предоставляет блочное (RADOS Block Device, RBD), файловое (CephFS) и объектное (RADOS Gateway, RGW) хранилище в единой системе \cite{ceph_arch}. Архитектура строится на нескольких типах демонов. \textit{Monitor (MON)} --- поддерживает главную копию cluster map и обеспечивает консенсус между узлами через алгоритм Paxos; для HA требуется не менее трёх мониторов. \textit{Manager (MGR)} --- отвечает за мониторинг, управление и расширения (плагины). \textit{OSD Daemon} --- хранит данные на физических дисках, выполняет репликацию и проверку состояния. \textit{Metadata Server (MDS)} --- управляет метаданными CephFS и обеспечивает POSIX-совместимый доступ к файловой системе.

Децентрализованное размещение данных реализовано через алгоритм CRUSH (Controlled Replication Under Scalable Hashing), который позволяет клиентам самостоятельно вычислять местоположение данных без обращения к центральному индексу. Это устраняет масштабируемые ограничения централизованного метаданных-сервера.

\predheading{CephFS в Kubernetes.} Rook автоматизирует создание и управление пулами CephFS и предоставляет CSI-драйвер, поддерживающий режим доступа ReadWriteMany (RWX). Это позволяет нескольким подам одновременно монтировать один том для доступа к общим данным --- актуальный сценарий для ML-инференса, где несколько реплик читают одну модель.

\predheading{Ограничения.} Развёртывание Ceph требует специализированных знаний: настройка пулов CRUSH, управление OSD-журналами, диагностика медленных OSD и оптимизация производительности под конкретную нагрузку создают значительный барьер для команд без опыта работы с Ceph. Минимальный production-кластер требует нескольких выделенных узлов хранения с отдельными дисками.

Ceph является универсальной production-системой хранения, а не специализированным решением для кэширования весов ML-моделей. Его возможности существенно превосходят необходимые для данной задачи, а операционные требования несоразмерны её масштабу. Для сценария бессерверного инференса, где ключевая метрика --- latency первого байта при холодном старте, --- Rook/Ceph обеспечивает избыточную функциональность при высокой стоимости эксплуатации.

\subsection{Сравнительная таблица}
\label{sec:comparison_table}

Для удобства сопоставления результаты анализа сведены в таблицу~\ref{tab:solutions_comparison}.

\begin{table}[htbp]
\centering
\caption{Сравнение решений для хранения и кэширования ML-моделей}
\label{tab:solutions_comparison}
\begingroup
\setlength{\tabcolsep}{1.5pt}
\renewcommand{\arraystretch}{1.15}
\renewcommand{\tabularxcolumn}[1]{m{#1}}
\tablebodyfont
\begin{tabularx}{\textwidth}{|>{\raggedright\arraybackslash}m{31mm}|*{5}{>{\centering\arraybackslash}X|}}
\hline
\textbf{Критерий} & \textbf{Alluxio} & \textbf{JuiceFS} & \textbf{Dragonfly} & \textbf{CubeFS} & \textbf{Rook/Ceph} \\ \hline
Лок.\ кэш на ноде        & \shortstack{Да\\(RAM/SSD)} & \shortstack{Да\\(диск)} & \shortstack{Да\\(диск)} & \shortstack{Да\\(диск)} & \shortstack{Да\\(OSD)} \\ \hline
P2P-доставка              & Нет & Нет & Да & Нет & Нет \\ \hline
FUSE/POSIX                & Частично & Да & Нет & Да & \shortstack{Да\\(CephFS)} \\ \hline
K8s-интеграция            & \shortstack{Helm/\\Operator} & CSI & DaemonSet & CSI & \shortstack{Rook\\Operator} \\ \hline
SPOF в архитектуре        & \shortstack{Master\\(OSS)} & \shortstack{Metadata\\Engine} & Scheduler & Master & \shortstack{MON\\(кворум)} \\ \hline
Объектное хранилище       & \shortstack{Опцио-\\нально} & \shortstack{Обяза-\\тельно} & \shortstack{Не\\нужно} & \shortstack{Не\\нужно} & \shortstack{Встроено\\(RGW)} \\ \hline
Сложность эксплуатации    & Высокая & Средняя & Средняя & Высокая & \shortstack{Очень\\высокая} \\ \hline
\shortstack[l]{Применимость\\к serverless} & Средняя & Низкая & Высокая & Средняя & Низкая \\ \hline
Лицензия                  & \shortstack{Community/\\Ent.} & \shortstack{Apache\\2.0} & \shortstack{Apache\\2.0} & \shortstack{Apache\\2.0} & \shortstack{LGPL/\\Apache} \\
\hline
\end{tabularx}
\endgroup
\end{table}

Анализ таблицы~\ref{tab:solutions_comparison} позволяет сформулировать следующие наблюдения. Ни одно из рассмотренных решений не удовлетворяет всем требованиям одновременно. Dragonfly наиболее близок по концепции P2P-доставки и показывает высокую применимость к serverless-сценариям, однако не предоставляет FUSE-интерфейса и требует отдельной инфраструктуры с собственным Scheduler. JuiceFS имеет POSIX-интерфейс и хорошую архитектуру метаданных, но обязательная зависимость от объектного хранилища и отсутствие P2P между нодами делают его неоптимальным для данной задачи. Alluxio и Rook/Ceph обладают богатой функциональностью, но их операционная сложность несоразмерна специализированной задаче кэширования весов. CubeFS занимает промежуточное положение: поддерживает POSIX и K8s-интеграцию, но лишена P2P-доставки и сложна в эксплуатации.

\subsection{Обоснование разработки специализированного решения}
\label{sec:justification}

Проведённый анализ демонстрирует системное несоответствие между возможностями существующих решений и требованиями задачи. Каждая из рассмотренных систем спроектирована как универсальная инфраструктура хранения данных общего назначения. Задача же данной работы принципиально уже: обеспечить быстрый локальный или межнодовой доступ к фрагментам (чанкам) файлов весов ML-моделей для контейнеров бессерверного инференса в Kubernetes. Ниже приведены конкретные аргументы в пользу разработки специализированного решения.

\predheading{Специализация вместо универсальности.} Все рассмотренные системы обеспечивают функциональность, существенно выходящую за рамки необходимой: поддержку множественных протоколов доступа, управление большими кластерами хранения, операции записи для приложений. Для задачи serverless-инференса нужен узкий cache/fetch слой, ориентированный на оптимизацию чтения крупных бинарных файлов по фрагментам. Дополнительная функциональность увеличивает поверхность отказов и затраты на эксплуатацию без соответствующего увеличения полезности.

\predheading{Приоритет локальности данных.} В идеальном случае чанк весов модели должен читаться с локального диска той же ноды, на которой запущен inference-контейнер. Это исключает сетевой хоп и обеспечивает минимальную задержку. Ни одна из рассмотренных систем не реализует явного приоритета <<локальный диск -> соседняя нода -> внешний источник>> в виде встроенного трёхуровневого read path для serverless-сценария.

\predheading{P2P в сочетании с FUSE.} Dragonfly реализует P2P-доставку, но не предоставляет POSIX-интерфейса. JuiceFS предоставляет FUSE, но не поддерживает P2P между нодами. Требуемая комбинация --- P2P-доставка чанков между узлами кластера с представлением результата через стандартный файловый путь --- отсутствует ни в одном из существующих решений в готовом виде.

\predheading{Простая модель метаданных с TTL.} Redis в роли реестра метаданных чанков обеспечивает атомарные операции, низкую задержку lookup ($O(1)$ по ключу) и встроенный механизм TTL для автоматического удаления устаревших записей о недоступных нодах. Это принципиально проще, чем поддержка полноценного Metadata Engine с журналами, репликой Raft и сложной схемой партиционирования, как в JuiceFS или CubeFS. Для кластеров малого и среднего размера (до тысяч нод) одиночный или сентинельный Redis является достаточным и значительно более управляемым решением.

\predheading{Fallback без жёсткой зависимости.} Специализированная система допускает ситуацию, когда нужный чанк отсутствует в кластере: в этом случае он скачивается из внешнего источника (Hugging Face или другого registry) и сохраняется в локальном кэше. Такой fallback требуется лишь при первом обращении к модели; все последующие поды получают чанк из локального или соседнего кэша. Это свойство устраняет принципиальную зависимость от объектного хранилища как обязательного звена в пути чтения, в отличие от JuiceFS.

\predheading{Отказоустойчивость без центрального файлового сервера.} NFS как общее хранилище, рассмотренное в разделе~3, решает задачу устранения повторных загрузок из интернета, но создаёт центральное узкое место: все чтения крупных файлов проходят через один файловый сервер. Специализированная система использует локальные диски всех worker-нод как распределённый кэш: при отказе одной ноды чанки, хранившиеся на ней, могут быть получены от других нод или из внешнего источника. Redis с TTL автоматически удаляет устаревшие записи об отказавшем агенте. Такая модель принципиально более устойчива к частичным отказам, чем централизованный NFS.

Дополнительным аргументом является экономическая доступность: все описанные компоненты (Redis, FUSE-библиотека, Go-реализация агента) распространяются под открытыми лицензиями. Alluxio Enterprise с ключевыми функциями кэширования и Rook/Ceph с его операционными требованиями несут значительные финансовые и инфраструктурные затраты, неприемлемые для небольших ML-платформ.

\subsection{Выводы по разделу}

Рассмотренные решения --- Alluxio, JuiceFS, Dragonfly, CubeFS и Rook/Ceph --- предоставляют ценные архитектурные идеи, нашедшие отражение в разрабатываемой системе: разделение данных и метаданных (JuiceFS), P2P-дистрибуция между узлами кластера (Dragonfly), предоставление POSIX-совместимого интерфейса через FUSE (JuiceFS, CubeFS, CephFS).

Вместе с тем ни одно из существующих решений не обеспечивает требуемую комбинацию свойств: трёхуровневый read path с приоритетом локального кэша, P2P-доставку чанков между нодами и прозрачный FUSE-интерфейс для ML-runtime без обязательной зависимости от внешнего объектного хранилища и тяжёлой операционной инфраструктуры.

Это обосновывает разработку специализированной системы: lightweight chunk-кэша на локальных дисках worker-нод, сопровождаемого Redis-реестром метаданных чанков и FUSE-клиентом, обеспечивающим runtime-контейнеру стандартный файловый интерфейс. Описание промежуточного этапа --- подхода на основе NFS и координации загрузок --- приведено в разделе~3. Основная архитектура разрабатываемой системы описана в разделе~4.
